\section{Metodologia}
\label{sec:metodologia}

A metodologia do trabalho será organizada em quatro etapas: caracterização do ambiente, definição da carga experimental, coleta de métricas e análise dos resultados.

Na caracterização do ambiente, serão registrados os componentes de hardware e software utilizados nos experimentos. Devem ser documentados processador, memória principal, armazenamento, GPU quando disponível, sistema operacional, bibliotecas, ferramentas de monitoramento e versões relevantes.

Na definição da carga experimental, será selecionada uma tarefa relacionada à IA generativa ou uma carga computacional representativa. Possíveis opções incluem inferência com modelo de linguagem de pequeno porte, execução de operação matricial/tensorial, benchmark sintético ou comparação entre execução com diferentes configurações de recursos.

Durante a coleta, serão medidas métricas como tempo de execução, uso de CPU, uso de memória, utilização de GPU, vazão de armazenamento e temperatura, conforme a disponibilidade das ferramentas do laboratório. Cada experimento deverá ser repetido quando possível, registrando variações e condições de execução.

Por fim, os dados serão organizados em tabelas e analisados com foco na identificação de gargalos. A análise buscará relacionar os resultados às características da infraestrutura e aos conceitos de HPC discutidos na fundamentação teórica.
